\documentclass[11pt,letterpaper]{article}
\usepackage[T1]{fontenc}
\usepackage[utf8]{inputenc}
\usepackage{lmodern}
\usepackage[margin=1in,headheight=15pt]{geometry}
\usepackage{amsmath,amssymb,amsthm,mathtools,mathrsfs}
\usepackage{microtype,booktabs,array,longtable,enumitem}
\usepackage[dvipsnames]{xcolor}
\usepackage{fancyhdr,aliascnt}
\usepackage[colorlinks=true,linkcolor=MidnightBlue,citecolor=MidnightBlue,urlcolor=MidnightBlue]{hyperref}
\usepackage[nameinlink,noabbrev]{cleveref}
\hypersetup{pdftitle={Holomorphic Dynamics on Surcomplex Halos: Exact Flows and Sharp Small-Divisor Thresholds},pdfauthor={Research manuscript prepared for Vladimir Reshetnikov},pdfsubject={Hahn-supported analytic dynamics, linearization, flows, resonances}}
\pagestyle{fancy}\fancyhf{}
\fancyhead[L]{\small Holomorphic dynamics on surcomplex halos}
\fancyhead[R]{\small\thepage}
\renewcommand{\headrulewidth}{0.3pt}
\setlength{\parskip}{0.25em}
\setlength{\parindent}{1.3em}
\setlength{\emergencystretch}{2em}
\setlist{nosep,leftmargin=2em}
\allowdisplaybreaks[2]
\newtheorem{theorem}{Theorem}[section]
\newaliascnt{lemma}{theorem}
\newtheorem{lemma}[lemma]{Lemma}
\aliascntresetthe{lemma}
\newaliascnt{proposition}{theorem}
\newtheorem{proposition}[proposition]{Proposition}
\aliascntresetthe{proposition}
\newaliascnt{corollary}{theorem}
\newtheorem{corollary}[corollary]{Corollary}
\aliascntresetthe{corollary}
\theoremstyle{definition}
\newaliascnt{definition}{theorem}
\newtheorem{definition}[definition]{Definition}
\aliascntresetthe{definition}
\newaliascnt{example}{theorem}
\newtheorem{example}[example]{Example}
\aliascntresetthe{example}
\newaliascnt{question}{theorem}
\newtheorem{question}[question]{Question}
\aliascntresetthe{question}
\theoremstyle{remark}
\newaliascnt{remark}{theorem}
\newtheorem{remark}[remark]{Remark}
\aliascntresetthe{remark}
\crefname{lemma}{lemma}{lemmas}
\Crefname{lemma}{Lemma}{Lemmas}
\crefname{proposition}{proposition}{propositions}
\Crefname{proposition}{Proposition}{Propositions}
\crefname{corollary}{corollary}{corollaries}
\Crefname{corollary}{Corollary}{Corollaries}
\crefname{definition}{definition}{definitions}
\Crefname{definition}{Definition}{Definitions}
\crefname{example}{example}{examples}
\Crefname{example}{Example}{Examples}
\crefname{question}{question}{questions}
\Crefname{question}{Question}{Questions}
\crefname{remark}{remark}{remarks}
\Crefname{remark}{Remark}{Remarks}
\newcommand{\C}{\mathbb C}
\newcommand{\R}{\mathbb R}
\newcommand{\Q}{\mathbb Q}
\newcommand{\Z}{\mathbb Z}
\newcommand{\N}{\mathbb N}
\newcommand{\No}{\mathbf{No}}
\newcommand{\SC}{\mathbf{SC}}
\newcommand{\HH}{\mathscr H}
\newcommand{\OO}{\mathcal O}
\newcommand{\GG}{\mathscr G}
\newcommand{\XX}{\mathscr X}
\newcommand{\m}{\mathfrak m}
\newcommand{\id}{\operatorname{id}}
\newcommand{\supp}{\operatorname{supp}}
\newcommand{\st}{\operatorname{st}}
\newcommand{\Cent}{\operatorname{Cent}}
\newcommand{\Fix}{\operatorname{Fix}}
\newcommand{\ord}{\operatorname{ord}}
\newcommand{\Span}{\operatorname{span}}
\newcommand{\DR}{\mathbb D_R^d}
\newcommand{\Li}{\mathcal L}
\newcommand{\NN}{\mathcal N}
\newcommand{\vmin}{v_{\min}}
\newcommand{\coef}[2]{[t^{#1}]#2}
\newcommand{\word}{\mathcal W}
\newcolumntype{P}[1]{>{\raggedright\arraybackslash}p{#1}}
\numberwithin{equation}{section}

\begin{document}
\hypersetup{pageanchor=false}
\begin{titlepage}
\centering
\vspace*{0.45in}
{\large\scshape A research contribution to surcomplex analysis\par}
\vspace{0.45in}
{\Huge\bfseries Holomorphic Dynamics\par}
\vspace{0.15in}
{\Huge\bfseries on Surcomplex Halos\par}
\vspace{0.35in}
{\Large Exact Flows and Sharp Small-Divisor Thresholds\par}
\vspace{0.35in}
{\large September 21, 2026\par}
\vspace{0.45in}
\begin{minipage}{0.92\textwidth}
\begin{abstract}
We study nonlinear maps with infinitesimal Hahn-supported perturbations over
set-sized fields $K_\Gamma=\C((t^\Gamma))\subset\No[i]$. The main result is
an exact classification of universal linearization for a nonresonant diagonal
unitary multiplier. If $\tau$ is the exponential growth rate of the inverse
homological divisors, then every positive perturbation is linearizable on its
\emph{unchanged ordinary coefficient polydisk} exactly when $\tau=0$.
For radius-free convergent germs, and separately for entire coefficients, the
criterion is $\tau<\infty$. Polynomial and formal coefficients require only
nonresonance. Necessity is witnessed at a single infinitesimal weight;
sufficiency follows from a support-controlled inverse of the nonlinear
Schr\"oder operator. We prove finite dependence at every Hahn exponent and
quantitative radius and degree bounds, without a rank-one or grid-support
assumption.

We also construct exact logarithms, fractional iterates, and flows of every
positive perturbation of the identity on the same coefficient domain. All
nonzero integer iterates have the same fixed-point ideal; in one variable the
full positive centralizer is an explicit group of Hahn times. A finite-order
multiplier admits a same-domain rotation--flow normal form. An explicit
non-Brjuno rotation satisfying $\tau=0$ separates the Hahn theory from ordinary
analytic linearization, and a critical-scale example explains why this is not
an all-scale theorem.
\end{abstract}
\end{minipage}
\vfill
\begin{minipage}{0.92\textwidth}\small
\textbf{Status and provenance.} This is an AI-assisted research manuscript,
not a refereed or formally verified paper. The proposed original contribution
is the precise arbitrary-support analytic package and its sharp coefficient-
category distinctions, not the classical formal logarithm or the classical
small-divisor problem. A targeted literature review found no exact published
match to the main package; it does not certify priority. No named published
conjecture is claimed solved. Twenty-one exact finite-jet checks accompany the
proofs but do not replace them. The motivating repository is reviewed at
commit \texttt{39f2be6667ade51bca2b45daa47e289d69c09764}.
\end{minipage}
\end{titlepage}
\hypersetup{pageanchor=true}
\pagenumbering{roman}
\tableofcontents
\clearpage
\pagenumbering{arabic}

\section{The question and the contribution}\label{sec:intro}

\subsection{From analytic functions to analytic dynamics}
The repository \texttt{VladimirReshetnikov/Surreal} develops Hahn-supported
surcomplex function theory, polynomial algebra, finite analytic geometry,
contour functionals, and global divisor constructions \cite{Repo}. Its
catalogue and its analytic-geometry report emphasize a distinction that is
particularly consequential for dynamics: holomorphic Hahn coefficients may
share a fixed ordinary domain, may merely define separate convergent germs,
or may be formal power series. These are different coefficient categories.
An argument proving a formal conjugacy does not automatically produce a
conjugacy in either analytic category.

Here is the dynamical question addressed in this paper. Suppose a holomorphic
map has an ordinary, explicitly understood reduction, but its higher data
contain arbitrarily many infinitesimal scales. Can one classify its dynamics
by a change of coordinates which has the \emph{same analytic coefficient
quality} as the original map? Even near a diagonal rotation, classical small
divisors make this a genuine issue. At the opposite extreme, when the
ordinary reduction is the identity, a formal logarithm is natural, but one
must still show that it is a strongly summable Hahn object on the original
coefficient domain, not a divergent expression or an unjustified topological
limit.

We answer these questions for positive-support perturbations. The decisive
resource is not a real norm bound. It is the finite number of ways to assemble
one Hahn exponent from a well-ordered positive support. That finiteness lets
us perform an arbitrarily deep nonlinear construction while doing only
finitely many ordinary analytic operations at each specified exponent.

\subsection{The sharp linearization statement}
Let $\Gamma$ be a nonzero set-sized ordered abelian group. For the surreal
interpretation we take it to be a divisible subgroup of $\No$. Write
$\HH_\Gamma(B)=B((t^\Gamma))$ for the Hahn algebra with coefficient algebra
$B$, and append a subscript $+$ for strictly positive support. Define
\[
 \DR=\{z\in\C^d: |z_j|<R\text{ for every }j\},\qquad 0<R<\infty.
\]
Fix a diagonal matrix $\Lambda=\operatorname{diag}(\lambda_1,\ldots,\lambda_d)$
with $|\lambda_j|=1$. Throughout the linearization theorem assume
\begin{equation}\label{eq:nonres}
 \lambda^\alpha\ne\lambda_j
 \quad (|\alpha|\ge2,\ 1\le j\le d),\qquad
 \lambda^\alpha=\lambda_1^{\alpha_1}\cdots\lambda_d^{\alpha_d}.
\end{equation}
Its small-divisor exponent is
\begin{equation}\label{eq:tauintro}
 \tau(\Lambda)=\limsup_{m\to\infty}\frac1m\log M_m(\Lambda),\qquad
 M_m(\Lambda)=\max\left(1,\max_{\substack{|\alpha|=m\\1\le j\le d}}
                    |\lambda^\alpha-\lambda_j|^{-1}\right).
\end{equation}
Thus $\tau\in[0,\infty]$. We consider
\begin{equation}\label{eq:Fintro}
 F(z)=\Lambda z+f(z),\qquad f\in\HH_\Gamma(B)_+^d,
 \qquad f_\gamma(0)=0,\quad df_\gamma(0)=0.
\end{equation}
The last condition is important: the multiplier is \emph{exactly} $\Lambda$,
not a matrix with an infinitesimal perturbation.

\begin{theorem}[Universal coefficient-category classification]\label{thm:main}
For every nonzero $\Gamma$ and every $\Lambda$ satisfying
\eqref{eq:nonres}, the following table gives necessary and sufficient
conditions for \emph{every} map \eqref{eq:Fintro} in the indicated category to
admit a normalized conjugacy
\begin{equation}\label{eq:Schroeder}
 H(z)=z+h(z),\qquad h\in\HH_\Gamma(B)_+^d,\quad
 h_\gamma(0)=0,\quad dh_\gamma(0)=0,\qquad H\circ F=\Lambda H.
\end{equation}
Whenever it exists, that conjugacy is unique among normalized positive-support
conjugacies, and its inverse belongs to the same category.
{\normalfont
\begin{center}
\begin{tabular}{@{}P{0.65\textwidth}P{0.25\textwidth}@{}}
\toprule
Coefficient category $B$ & Exact condition\\
\midrule
$\OO(\DR)$, on the \emph{same fixed} $\DR$ & $\tau(\Lambda)=0$\\
$\C\{z_1,\ldots,z_d\}$, with no common radius required & $\tau(\Lambda)<\infty$\\
$\OO(\C^d)$, entire coefficients & $\tau(\Lambda)<\infty$\\
$\C[z_1,\ldots,z_d]$, polynomial coefficients & Nonresonance only\\
$\C[[z_1,\ldots,z_d]]$, formal coefficients & Nonresonance only\\
\bottomrule
\end{tabular}
\end{center}}
If $S=\supp f\subset\Gamma_{>0}$, then $\supp h$ is contained in the
positive part of the additive monoid generated by $S$. Each $h_\gamma$ depends
on only finitely many coefficients of $f$ and finitely many ordinary
homological inversions. The result is compatible with enlargement of the
Hahn workspace.
\end{theorem}
The proof occupies \cref{sec:homological,sec:linearization,sec:quantitative}.
Necessity is already visible in an input with support $\{\eta\}$, for any
$\eta>0$. Sufficiency is genuinely a nonlinear arbitrary-support statement:
no countability, finite rank, or finitely generated grid assumption is made
on $S$.

\subsection{Exact iteration and resonance}
For an arbitrary open $U\subset\C^d$, every map $F=\id+u$ with
$u\in\HH_\Gamma(\OO(U))_+^d$ has a unique positive Hahn vector field $X$
whose time-one map is $F$ (\cref{thm:exp-log}). All its ordinary complex-time
maps have holomorphic coefficients on that \emph{same} $U$. No shrinking and
no uniform bounds on the coefficient functions are required.

This yields more than fractional iteration. The fixed-point ideals of all
nonzero integer iterates are identical, so neither a new periodic point nor
a new fixed-point multiplicity can appear by iteration
(\cref{thm:fixed-ideal}). In one variable on a connected ordinary domain, the
entire centralizer within the positive substitution group consists of
explicit admissible Hahn times of this vector field
(\cref{thm:centralizer}). When the ordinary multiplier has finite order, a
finite averaging construction produces a commuting rotation and resonant
positive flow on the original polydisk (\cref{thm:finite-order}).

\subsection{What is, and is not, new}
Formal logarithms and exponentials of tangent-to-identity maps, their
relation to vector fields, and semisimple--unipotent decompositions are
classical. Rib\'on provides a close modern treatment in the formal and smooth
settings \cite{Ribon}. Small-divisor linearization is also classical; we use
the work of Carletti--Marmi and the Brjuno--Yoccoz theory to identify the
relevant predecessor and to compare the arithmetic, not to relabel those
theorems as new \cite{CarlettiMarmi,Yoccoz,BuffCheritat}.

The proposed contribution here is the precise conjunction of arbitrary
well-ordered Hahn support, preservation or controlled loss of ordinary
coefficient domains, sharp universal necessity, finite dependence at a
specified exponent, and the resulting dynamical statements on surcomplex
halos. The main arithmetic distinction is meaningful: there are explicitly
constructed non-Brjuno rotations with $\tau=0$, and they satisfy the first
row of \cref{thm:main} even though the corresponding ordinary quadratic map
is not analytically linearizable. We prove the construction and explain the
absence of a contradiction in \cref{sec:arithmetic}.

\section{Workspaces, coefficient algebras, and evaluation}\label{sec:framework}

\subsection{A set-sized setting inside the surcomplex numbers}
A Hahn series with coefficients in a commutative $\C$-algebra $B$ is a formal
sum
\[
 A=\sum_{\gamma\in\Gamma}t^\gamma A_\gamma,
 \qquad A_\gamma\in B,
\]
whose support $\{\gamma:A_\gamma\ne0\}$ is well ordered in the given order
on $\Gamma$. We use the standard convolution product. For a nonzero series
put $v(A)=\min\supp A$, and let $v(0)=\infty$. For a vector or matrix, take
the minimum valuation of its entries. The condition $v(A)>0$ means that
\emph{all} of its nonzero coefficients have positive exponents.

When $B=\C$, this is the field $K_\Gamma=\C((t^\Gamma))$. Its valuation
ring and maximal ideal are
\[
 \OO_{K_\Gamma}=\{x:v(x)\ge0\},\qquad
 \m_{K_\Gamma}=\{x:v(x)>0\}.
\]
We regard $\C$ as its constant coefficient field. If $\Gamma\subset\No$,
the substitution $t^\gamma=\omega^{-\gamma}$ identifies these series with
surcomplex normal forms. All objects used in any one construction are sets.
There is no ring of proper-class-supported series in the argument.

For completeness, any \emph{set} of surcomplex constants can be included in
one divisible workspace: collect the exponents in the real and imaginary
Conway normal forms of those constants, change their signs to our $t$
convention, and take their $\Q$-linear span in $\No$. The union of the
supports is a set, and so is this span. Later constructions use sums of
exponents already in that group. Evaluation at additional surcomplex points
is handled by enlarging the group to include their supports. Thus the
workspace convention is a size discipline, not a restriction to one
Archimedean scale. The broader surreal analytic setting and its composition
subtleties are discussed by Berarducci--Mantova \cite{BM}; our composition is
the explicit coordinate Taylor composition defined below.

\begin{definition}[Strong summability]\label{def:summable}
A set-indexed family $(A_i)_{i\in I}$ of Hahn series is \emph{strongly
summable} if $\bigcup_i\supp A_i$ is well ordered and each exponent belongs
to the support of only finitely many members. Its sum is then defined by the
finite coefficient sums. Every infinite sum in the Hahn direction below
uses this convention.
\end{definition}
Ordinary Taylor series of coefficient functions, on the other hand, use
ordinary complex convergence. Keeping these two summations separate is
essential. In particular, we never justify a Hahn sum by taking a limit of
partial sums in the fine order topology of $\No$.

\subsection{Five categories, and a sixth question we leave distinct}
We will use the coefficient algebras
\begin{equation}\label{eq:categories}
 \OO(U),\qquad \C\{z\},\qquad \OO(\C^d),\qquad
 \C[z],\qquad \C[[z]],
\end{equation}
with vector-valued versions understood componentwise. The notation $\C\{z\}$
means convergent power-series germs at the origin. In
$\HH_\Gamma(\C\{z\})$, each $A_\gamma$ has some positive radius, but the
radii need not have a common positive lower bound. In
$\HH_\Gamma(\OO(\DR))$, the entire coefficient family is holomorphic on
the same specified ordinary polydisk, with no uniform norm bound imposed.

A common-domain \emph{germ} algebra, formed by allowing one common polydisk
and then shrinking it, is yet another object. The first row of
\cref{thm:main} concerns preservation of a \emph{specified unchanged}
polydisk. The second row concerns separate convergent germs. Neither row,
nor their combination, classifies all common-domain germs in the intermediate
case $0<\tau<\infty$. We return to this genuinely separate question in
\cref{sec:limitations}.

The strict distinction is easy to witness. For $\eta>0$,
\begin{equation}\label{eq:shrinking-radii}
 \sum_{n\ge1}\frac{t^{n\eta}}{1-nz}
\end{equation}
has convergent-germ coefficients but no common ordinary disk. By contrast,
$\sum_{n\ge0}n!z^n$ is a formal coefficient with zero complex radius. Neither
its formal status nor the lack of a common radius in \eqref{eq:shrinking-radii}
prevents evaluation at an infinitesimal argument in the sense below.

\subsection{Halos and monads}
For an ordinary domain $U\subset\C^d$, set
\[
 U_\Gamma^\#=\{a+\xi:a\in U,\ \xi\in\m_{K_\Gamma}^d\}.
\]
The constant coefficient $a$ is its standard part. Given
$A=\sum_\gamma t^\gamma A_\gamma\in\HH_\Gamma(\OO(U))$, define
\begin{equation}\label{eq:halo-evaluation}
 A(a+\xi)=
 \sum_{\gamma,\alpha}t^\gamma
 \frac{\partial^\alpha A_\gamma(a)}{\alpha!}\xi^\alpha.
\end{equation}
The family is strongly summable by \cref{lem:neumann,lem:shifted} below,
applied to the finite union of the positive supports of the components of
$\xi$. For a germ or formal coefficient, use its Taylor coefficients at
$0$ and restrict the domain to $\m_{K_\Gamma}^d$. Polynomial and entire
coefficients also give evaluations on every finite halo, but not automatically
at every infinite surcomplex point.

At a fixed ordinary Taylor center, coefficients such as $n!$ or extremely
large inverse small divisors remain \emph{ordinary complex constants}.
Their valuation is zero. Strong summability does not impose a real growth
condition on them. This fact explains both the strength of the forthcoming
theorems and the need to state their domains exactly.

\begin{remark}[The derivative in this paper]
The symbol $\partial_j$ differentiates the ordinary coordinate $z_j$ and
acts coefficientwise; every $t^\gamma$ is constant for this derivative.
It is not a derivation of the surreal coefficient field such as the
Berarducci--Mantova derivation. No compatibility with a global composition
operator on all of $\No$ is asserted.
\end{remark}

\section{The support certificate: finite words, not sequential limits}
\label{sec:support}

\subsection{The positive-word lemma}
Let $S\subset\Gamma_{>0}$ be well ordered. Let $S^*$ denote its additive
monoid, including $0$. For $\gamma\in S^*$, define
\[
 \word_S(\gamma)=
 \{(s_1,\ldots,s_k): k\ge0,\ s_i\in S,\ s_1+\cdots+s_k=\gamma\}.
\]
Here words are ordered, so permutations are counted separately.

\begin{lemma}[Neumann's support lemma, in finite-word form]\label{lem:neumann}
The set $S^*$ is well ordered, and $\word_S(\gamma)$ is finite for every
$\gamma\in\Gamma$. In particular,
\begin{equation}\label{eq:depth}
 \ell_S(\gamma)=\max\{k:(s_1,\ldots,s_k)\in\word_S(\gamma)\}
\end{equation}
is finite whenever $\gamma\in S^*$, with $\ell_S(0)=0$.
\end{lemma}
\begin{proof}
We give the standard well-quasi-order argument, making the finite-dependence
conclusion explicit. Order finite words over $S$ by subsequence embedding
with coordinatewise domination. Higman's lemma says that this word order is
a well-quasi-order: every infinite sequence contains indices $i<j$ for which
the earlier word embeds in the later one \cite{Higman}. Because all letters
are positive, such an embedding weakly increases the sum of the letters.
It increases that sum strictly unless the two words are identical.

An infinite strictly decreasing sequence in $S^*$ would, after choosing one
representing word for each element, contradict the weak increase just
observed. Thus $S^*$ is well ordered. Infinitely many distinct words of one
fixed weight would contradict the strict-increase observation. This proves
both assertions.
\end{proof}
These are classical support facts, not new claims. The invertibility and
support viewpoint is also developed in Poonen's recent account of
Hahn--Mal'cev--Neumann rings \cite{Poonen}. A short proof of the form of
Higman's lemma just used is included in \cref{app:higman}.

For $\gamma\in S^*$ let $D_S(\gamma)$ be the set of letters occurring in
all words in $\word_S(\gamma)$. It is finite. Finite partial sums of these
words will serve as intermediate weights. This is the certificate that
permits a computation at weight $\gamma$ without traversing every smaller
Hahn exponent.

\begin{lemma}[Adding an arbitrary input support]\label{lem:shifted}
If $A\subset\Gamma$ is well ordered, then $A+S^*$ is well ordered. For any
fixed $\gamma$, there are only finitely many pairs consisting of $a\in A$
and a word in $S$ whose total weight, after addition of $a$, is $\gamma$.
\end{lemma}
\begin{proof}
For nonempty $A$, subtract $a_0=\min A$. The finite union
\[
 T=S\cup\bigl((A-a_0)\setminus\{0\}\bigr)
\]
is a well-ordered positive set. Apply \cref{lem:neumann} to $T$ and to
weight $\gamma-a_0$. A pair in the assertion is encoded by a word with one
distinguished initial letter from $A-a_0$, or no such letter if $a=a_0$,
followed by its $S$-word. There are only finitely many encodings at that
weight. The well-ordering assertion follows from $A+S^*\subset a_0+T^*$.
\end{proof}

\subsection{Support-increasing operators}
Suppose $T$ is a $\C$-linear operator on coefficient families such that every
term of $T(A)$ adds a nonempty word from $S$ to an input exponent of $A$,
with finitely many coefficient operations for each such word. Then
\begin{equation}\label{eq:operator-series}
 \sum_{k\ge0}c_kT^k(A)
\end{equation}
is strongly summable for any ordinary scalars $c_k$. At each output exponent,
only finitely many input exponents, word lengths, and operations contribute.
Consequently polynomial identities for a nilpotent operator may be used
coefficientwise in \eqref{eq:operator-series}. This justifies geometric,
logarithmic, and exponential identities below.

This principle also applies to operators with coefficients in a Hahn algebra
and to vector or matrix operators. It applies to a strongly summable input
family simultaneously: \cref{lem:shifted} bounds the relevant input weights,
and local finiteness of that family bounds the number of input members.
Thus the operators we construct commute with legitimate strong regrouping.

\begin{example}[Infinitely many earlier weights, but only one relevant word]
\label{ex:ranktwo}
Take $\Gamma=\Q+\Q\Omega$, ordered with $\Omega$ greater than every rational,
and
\[
 S=\{1\}\cup\{2-1/n:n\ge2\}\cup\{\Omega\}.
\]
This is well ordered, but is not contained in a finitely generated additive
grid. There are infinitely many members of $S$ below $\Omega$. Nevertheless,
\[
 \word_S(\Omega)=\{(\Omega)\},\qquad
 \ell_S(\Omega)=1,\qquad \ell_S(\Omega+2)=3.
\]
The words at $\Omega+2$ are the three permutations of $(\Omega,1,1)$.
Any finite sum of rational letters is rational, and any sum involving two
$\Omega$'s is larger than $\Omega+2$. These observations prove the claims.
An algorithm computing the coefficient at $\Omega$ needs the input at
$\Omega$, not infinitely many rational predecessors.
\end{example}

\begin{remark}[Why a usual contraction argument would be incomplete]
Writing $v(T^kA)\ge v(A)+k\min S$ does not imply that these valuations
become cofinal in $\Gamma$. For example, $k<\Omega$ for every integer $k$.
The finite-word argument establishes summability even when that sequence of
lower bounds never passes a specified higher rank. Every use of an infinite
operator series in this paper rests on this stronger certificate.
\end{remark}

\section{Positive substitution on an unchanged ordinary domain}
\label{sec:substitution}

\subsection{The substitution algebra}
For $U\subset\C^d$ open, define
\[
 \GG_+(U)=\{F(z)=z+u(z):u\in\HH_\Gamma(\OO(U))_+^d\}.
\]
No condition at $z=0$ is imposed in this definition, and $0$ need not belong
to $U$. If $F=\id+u$, its pullback on $A\in\HH_\Gamma(\OO(U))$ is
\begin{equation}\label{eq:pullback}
 C_F A=A\circ F=
 \sum_{\alpha\in\N^d}\frac{1}{\alpha!}
                 (\partial^\alpha A)u^\alpha.
\end{equation}
Every coefficient of this expression is a finite sum of products and
derivatives of coefficient functions holomorphic on $U$. Thus it remains
holomorphic on $U$, regardless of the size of its ordinary complex norm.
The support is contained in $\supp A+S^*$, where $S$ is the union of the
supports of the components of $u$.

\begin{proposition}[Substitution group]\label{prop:subgroup}
The rule \eqref{eq:pullback} defines an associative group law on
$\GG_+(U)$. Every inverse again belongs to $\GG_+(U)$, on the same $U$.
Evaluation gives mutually inverse maps of $U_\Gamma^\#$ and commutes with
composition. The same statements hold in the germ, entire, polynomial, and
formal coefficient categories, with the evaluation domains from
\cref{sec:framework}.
\end{proposition}
\begin{proof}
Strong summability follows from \cref{lem:neumann,lem:shifted}. At a fixed
Hahn exponent, only finitely many derivatives and products contribute.
Introduce a formal parameter for each of the finitely many input support
letters involved, and discard parameter monomials beyond the required word
length. Taylor composition in this finite nilpotent parameter algebra is
associative and multiplicative. Each coefficient identity in the asserted
group law is one of these identities. Equivalently, ordinary formal Taylor
identities may be expanded and regrouped using the same finiteness proof.

For the inverse, put $\Delta=C_F-I$. Every term of $\Delta$ adds a nonempty
word from $S$. Hence
\[
 C_F^{-1}=\sum_{k\ge0}(-\Delta)^k
\]
is well defined coefficientwise and is a two-sided inverse as a linear
operator. Since $C_F$ is multiplicative, this inverse is multiplicative too.
Let $G=C_F^{-1}z$. It is identity plus positive support, with support in
$S^*$. On coordinate polynomials, multiplicativity gives $C_F^{-1}=C_G$.
At each Hahn weight both operators are finite-order differential operators
on an ordinary input function; agreement on polynomial jets makes them agree
on holomorphic functions as well. For formal coefficients, the corresponding
finite-jet argument is coefficientwise in $z$. Strong linearity then extends
the equality to Hahn inputs. Thus $C_G$ is the two-sided inverse of $C_F$,
and $G$ is the two-sided substitution inverse of $F$. The finite coefficient
operations preserve every stated coefficient category.

Finally, substitute $z=a+\xi$ in each identity. The combined positive support
set now also includes the supports of the components of $\xi$. The same
lemmas allow every regrouping, so Taylor evaluation preserves the identities.
The standard part of $F(a+\xi)$ is $a$, and hence the evaluated map stays in
$U_\Gamma^\#$.
\end{proof}

\begin{remark}[Pullback order]
With our convention,
\begin{equation}\label{eq:pullback-order}
 C_FC_G=C_{G\circ F}.
\end{equation}
Pullback reverses composition. Iterates of a single map are unaffected by
this reversal, and commuting maps give commuting pullbacks. We retain the
convention explicitly to avoid a sign or order ambiguity in logarithms.
\end{remark}

For maps of the form $F=\Lambda z+f$ on $\DR$, with diagonal unitary
$\Lambda$ and $f$ positive, use
\begin{equation}\label{eq:rotation-substitution}
 A\circ F=
 \sum_{\alpha\in\N^d}\frac{1}{\alpha!}
           (\partial^\alpha A)(\Lambda z) f^\alpha.
\end{equation}
The ordinary center remains in $\DR$, since $\Lambda\DR=\DR$.
This provides the composition used in the linearization and finite-order
theorems. For formal coefficients or germs, $F(0)=0$ is imposed whenever
composition with its nonidentity reduction is used.

\section{Exact logarithms, flows, and fractional iterates}\label{sec:flows}

\subsection{Logarithms are actual coefficientwise vector fields}
A positive Hahn vector field on $U$ is an operator
\[
 X=a_1(z)\partial_1+\cdots+a_d(z)\partial_d,
 \qquad a\in\HH_\Gamma(\OO(U))_+^d.
\]
Its exponential and the logarithm of a positive substitution are defined by
\begin{equation}\label{eq:explog}
 \exp X=\sum_{k\ge0}\frac{X^k}{k!},\qquad
 \log C_F=\sum_{k\ge1}\frac{(-1)^{k+1}}{k}(C_F-I)^k.
\end{equation}
These are operators on the Hahn algebra, not numerical logarithms of the
values of $F$.

\begin{theorem}[Same-domain exponential--logarithm correspondence]
\label{thm:exp-log}
For every open $U\subset\C^d$, the assignments
\begin{equation}\label{eq:maplog}
 F\longmapsto X_F=\log C_F,
 \qquad
 X\longmapsto\Phi_X=(\exp X)(z)
\end{equation}
are inverse bijections between $\GG_+(U)$ and positive Hahn vector fields
on $U$. All coefficient functions stay holomorphic on the original $U$.
If $S$ contains the positive support of the input, the output support is
contained in $S^*\setminus\{0\}$, and each output coefficient has a finite
word certificate. Moreover,
\[
 v(X_F)=v(F-\id)
\]
for $F\ne\id$. The same assertions hold with germ, entire, polynomial, or
formal coefficients.
\end{theorem}
\begin{proof}
\emph{Summability.} Put $\Delta=C_F-I$ and let $S=\supp(F-\id)$, taking
the union of the vector supports. Formula \eqref{eq:pullback} shows that
$\Delta$ adds nonempty words from $S$. The logarithmic series in
\eqref{eq:explog} is therefore strongly summable on every input, by
\cref{sec:support}. Every output coefficient is a finite differential
expression with holomorphic coefficients on $U$.

\emph{The derivation identity.} Introduce an ordinary indeterminate $s$ and
write
\[
 C_F^s=\sum_{k\ge0}\binom{s}{k}\Delta^k.
\]
Each Hahn coefficient of $C_F^s(A)$ is a polynomial in $s$. For every
nonnegative integer $s$, this is the usual integer power of the
multiplicative operator $C_F$. Hence the coefficientwise polynomial identity
\[
 C_F^s(AB)=C_F^s(A)C_F^s(B)
\]
holds for all $s$. Differentiating at $s=0$ gives
\[
 X_F(AB)=X_F(A)B+A X_F(B),
\]
since $\frac{d}{ds}\binom{s}{k}|_{s=0}=(-1)^{k+1}/k$ for $k\ge1$.
All operations are justified coefficientwise by finite polynomials.

\emph{Why the derivation is a vector field.} At each Hahn weight $\gamma$,
the restriction of $X_F$ to ordinary coefficient functions is a finite-order
differential operator $D_\gamma$. The preceding identity makes
$D_\gamma$ a derivation. Set $a_{j,\gamma}=D_\gamma z_j$. By Leibniz,
$D_\gamma$ agrees with $\sum_j a_{j,\gamma}\partial_j$ on every coordinate
polynomial. At any ordinary point, a finite-order differential operator is
determined by a finite Taylor jet, so the two operators agree on every
holomorphic function. For convergent or formal germs the same argument uses
finite jets at the origin; for polynomials it is immediate. Thus
$X_F=\sum_j a_j\partial_j$ with positive Hahn coefficients $a_j$.
This finite-order observation avoids any appeal to a classification of
arbitrary algebraic derivations of a function ring.

\emph{Conversely.} If $X=a\cdot\partial$ is positive, each application adds
a positive support letter. The exponential is strongly summable. Repeated
Leibniz gives
\[
 X^k(AB)=\sum_{j=0}^k\binom{k}{j}X^j(A)X^{k-j}(B),
\]
so $\exp X$ is multiplicative and has inverse $\exp(-X)$. Put
$\Phi_X=(\exp X)z$. It is identity plus positive support. On coordinate
polynomials, multiplicativity gives $\exp X=C_{\Phi_X}$. At each Hahn
weight both sides are finite-order differential operators on an ordinary
input function, so the finite-jet argument extends the equality to the
whole coefficient algebra and then to Hahn inputs.

\emph{Inverse identities and support.} The identities
$\exp(\log(I+T))=I+T$ and $\log(\exp T)=T$ hold in a nilpotent operator
algebra. At any specified Hahn coefficient their application here involves
only finitely many words, so they hold coefficientwise. This proves the
bijection. Finally, $X_Fz=(F-z)+\text{terms with at least two support
letters}$. If $\alpha=v(F-z)>0$, the latter terms have valuation at least
$2\alpha>\alpha$. Thus the first coefficient is unchanged and the valuation
claim follows. Derivatives and finite products preserve all stated
coefficient categories.
\end{proof}

\subsection{Time parameters and unique roots}
\begin{corollary}[Exact flow law]\label{cor:flow}
Let $X=a\cdot\partial$ be positive. For every ordinary $s\in\C$,
\[
 \Phi_s=(\exp(sX))z
\]
belongs to $\GG_+(U)$ and satisfies
\begin{equation}\label{eq:flow-law}
 \Phi_0=\id,\qquad
 \Phi_{s+r}=\Phi_s\circ\Phi_r,\qquad
 \frac{d}{ds}\Phi_s=a\circ\Phi_s.
\end{equation}
For a fixed support exponent, the coefficient of $\Phi_s$ is a polynomial in
$s$. If $c\in K_\Gamma$ and either $c=0$ or
\begin{equation}\label{eq:admissible-time}
 v(c)+v(a)>0,
\end{equation}
then $\Phi_c=(\exp(cX))z$ is also defined and positive. The additive flow
law holds for any such times.
\end{corollary}
\begin{proof}
All ordinary scalar coefficients have valuation zero. The series for each
coefficient involves finitely many powers of $s$, and commuting exponentials
of multiples of $X$ satisfy the addition law. Differentiating coefficientwise
and commuting $X$ with its exponential gives
\[
 \frac{d}{ds}(\exp(sX))z=(\exp(sX))(Xz)
                         =a\circ\Phi_s.
\]
For Hahn time $c$, condition \eqref{eq:admissible-time} says exactly that
$cX$ has positive support, since $c$ is a scalar independent of $z$.
Apply \cref{thm:exp-log} to $cX$. The positive supports of two admissible
multiples have well-ordered union, so the same word argument proves their
addition identity.
\end{proof}
The condition \eqref{eq:admissible-time} characterizes when this scalar
multiple of a nonzero vector field is positive. It is not asserted to be a
maximal possible lifetime in a larger class of maps. For example, a constant
vector field can have translations beyond the positive subgroup.
Likewise, \eqref{eq:flow-law} describes coefficientwise analytic time
dependence, not a continuous path in the fine surreal topology.

\begin{corollary}[Unique fractional iteration]\label{cor:roots}
For every $F\in\GG_+(U)$ and every integer $N\ge1$, there is a unique
$G\in\GG_+(U)$ with $G^{\circ N}=F$, namely
\[
 G=(\exp(X_F/N))z.
\]
The group $\GG_+(U)$ is uniquely divisible and torsion-free.
\end{corollary}
\begin{proof}
Existence follows from the flow law. If $G^{\circ N}=F$, then
$C_G^N=C_F$. By \cref{thm:exp-log}, taking logarithms gives
$N X_G=X_F$, whence uniqueness. If $F^{\circ N}=\id$, then $NX_F=0$;
characteristic zero gives $F=\id$.
\end{proof}

\section{Fixed-point schemes and one-variable centralizers}
\label{sec:fixed-central}

\subsection{All nonzero iterates have the same fixed-point ideal}
For a coefficient map $F$ let
\[
 I_{\Fix(F)}=(F_1-z_1,\ldots,F_d-z_d)
\]
be its fixed-point ideal in the chosen Hahn coefficient algebra. No
Noetherianity or Nullstellensatz is needed to compare such ideals.

\begin{theorem}[Fixed-point ideal of a positive flow]\label{thm:fixed-ideal}
Let $X=a\cdot\partial$ be a positive Hahn vector field on $U$. If $c\ne0$
is an admissible Hahn time, then
\begin{equation}\label{eq:fixed-ideal}
 I_{\Fix(\Phi_c)}=(a_1,\ldots,a_d)
 \quad\text{in }\HH_\Gamma(\OO(U)).
\end{equation}
The same equality holds in the other coefficient categories and after
localization. In particular, all nonzero integer iterates of
$F\in\GG_+(U)$ have identical fixed-point ideals. On the evaluated halo,
there are no nonfixed points of finite period. Whenever a fixed-point
quotient or local algebra has finite length, that length is unchanged by
nonzero integer iteration.
\end{theorem}
\begin{proof}
Set $E=cX=b\cdot\partial$, so $b=ca$ is positive. Write the coordinate
difference as
\[
 \Phi_c-z=b+\sum_{k\ge2}\frac{E^kz}{k!}.
\]
Since $E^kz_j=\sum_\ell b_\ell\partial_\ell(E^{k-1}z_j)$, define
\begin{equation}\label{eq:fixedmatrix}
 B_{j\ell}=\delta_{j\ell}+
       \sum_{k\ge2}\frac{\partial_\ell(E^{k-1}z_j)}{k!}.
\end{equation}
The support lemma gives strong summability, and $B=I+\text{positive}$.
Consequently $B$ is invertible by its matrix geometric series and
\begin{equation}\label{eq:fixedfactor}
 \Phi_c-z=Bb.
\end{equation}
The two coordinate ideals are therefore equal. Because the scalar $c$ is a
unit in $K_\Gamma$, the ideal generated by $b$ is the one generated by $a$.
This proves \eqref{eq:fixed-ideal} and its localized versions.

On evaluation at any halo point, $B$ remains invertible; its inverse
identity can be evaluated by \cref{prop:subgroup}. Thus a point is fixed by
$\Phi_c$ exactly when all components of $a$ vanish there. Apply this to
$c=1$ and any nonzero integer $c=n$ to exclude nonfixed periodic points.
Identical ideals give identical quotient algebras, so any finite length or
local multiplicity defined by them is unchanged.
\end{proof}

The theorem concerns maps whose ordinary reduction is identity. It makes no
claim that a general surcomplex map lacks periodic points. In particular,
rotations of finite order lie outside this positive group unless they are
identity.

\subsection{The full positive centralizer in one variable}
\begin{theorem}[Centralizer as an explicit Hahn-time group]
\label{thm:centralizer}
Let $U\subset\C$ be connected, let $F\in\GG_+(U)$ be nonidentity, and
write $X_F=a(z)\partial_z$, $\alpha=v(a)>0$. Then
\begin{equation}\label{eq:centralizer}
 \Cent_{\GG_+(U)}(F)=
 \bigl\{(\exp(ca\partial_z))z:
        c\in K_\Gamma,\ c=0\text{ or }v(c)>-\alpha\bigr\}.
\end{equation}
The parametrization is an isomorphism from the displayed additive group of
Hahn times to the centralizer. Zeros of $a$ do not require their removal
from $U$.
\end{theorem}
\begin{proof}
Let $G$ be positive and write $X_G=b\partial_z$. If $F$ and $G$ commute,
their pullbacks commute. Their logarithms, being coefficientwise sums of
powers of $C_F-I$ and $C_G-I$, commute as well. All double regroupings are
valid using the union of the two positive support sets. Conversely, commuting
logarithms have commuting exponentials. Thus
\[
 FG=GF\quad\Longleftrightarrow\quad
 [a\partial_z,b\partial_z]=0
 \quad\Longleftrightarrow\quad ab'-ba'=0.
\]
Embed the holomorphic coefficient algebra into
$\mathcal M(U)((t^\Gamma))$, where $\mathcal M(U)$ is the field of
meromorphic functions on the connected domain $U$. The nonzero series $a$
is invertible there: factor its leading term and use a positive geometric
series for the remaining factor. Hence
\[
 (b/a)'=(ab'-ba')/a^2=0.
\]
Every Hahn coefficient of $b/a$ is a meromorphic function with zero
derivative. Connectedness makes it an ordinary complex constant. Therefore
$b/a=c\in K_\Gamma$.

Conversely, $b=ca$ commutes with $a$ for every scalar $c$. It has positive
support exactly when $c=0$ or $v(c)+\alpha>0$. The exponential--logarithm
bijection proves \eqref{eq:centralizer}. If two times give the same map,
their logarithms give $(c_1-c_2)a=0$, and $a\ne0$ implies $c_1=c_2$.
The flow addition law proves the group assertion. Meromorphic division was
used only to identify the ratio; the resulting coefficient $ca$ is
holomorphic on the original $U$, including the zeros of $a$.
\end{proof}

\begin{remark}
The centralizer pattern is familiar in formal one-dimensional dynamics;
what is retained here is the full original coefficient domain and an explicit
valuation range for possibly infinite Hahn times. The connectedness
hypothesis is necessary to obtain a single scalar series, rather than
independent scalar series on different ordinary components. In dimension
$d>1$, the correct general statement is that commuting positive maps have
commuting positive logarithmic vector fields; those vector fields need not
be scalar multiples.
\end{remark}

\section{Homological divisors and the radius of one inversion}
\label{sec:homological}

\subsection{The diagonal homological operator}
Return to the nonresonant diagonal unitary multiplier $\Lambda$ in
\eqref{eq:nonres}. Let $B_2^d$ denote vectors of coefficients whose Taylor
series have no terms of total degree $0$ or $1$. Define
\begin{equation}\label{eq:homological}
 \Li h=h(\Lambda z)-\Lambda h(z).
\end{equation}
If $h_j(z)=\sum_{|\alpha|\ge2}h_{j,\alpha}z^\alpha$, then
\[
 (\Li h)_{j,\alpha}=(\lambda^\alpha-\lambda_j)h_{j,\alpha}.
\]
Thus nonresonance gives a unique inverse on formal coefficients:
\begin{equation}\label{eq:Linverse}
 (\Li^{-1}g)_j(z)=
 \sum_{|\alpha|\ge2}\frac{g_{j,\alpha}}
                         {\lambda^\alpha-\lambda_j}z^\alpha.
\end{equation}
It preserves total polynomial degree and acts coefficientwise on Hahn
families without changing their Hahn exponents.

For a vector formal power series $g$, define its isotropic Taylor radius by
\begin{equation}\label{eq:radius}
 \rho(g)^{-1}=
 \limsup_{m\to\infty}
       \left(\max_{\substack{|\alpha|=m\\j}}|g_{j,\alpha}|\right)^{1/m},
\end{equation}
with the usual conventions for $0$ and $\infty$. This is exactly the
radius of absolute convergence on centered equal-radius polydisks. Indeed,
the number of multiindices of total degree $m$ is polynomial in $m$, which
does not affect the root test. A holomorphic function on $\DR$ has radius
at least $R$, and an entire function has infinite radius.

\begin{lemma}[Single-step analytic loss]\label{lem:radius-loss}
If $\tau(\Lambda)<\infty$, then
\begin{equation}\label{eq:loss}
 \rho(\Li^{-1}g)\ge e^{-\tau(\Lambda)}\rho(g).
\end{equation}
In particular, $\Li^{-1}$ preserves $\OO(\DR)_2^d$ when $\tau=0$,
preserves convergent germs when $\tau<\infty$, and preserves entire
coefficients when $\tau<\infty$. It always preserves polynomial and formal
coefficients under nonresonance.
\end{lemma}
\begin{proof}
By \eqref{eq:Linverse}, the maximum degree-$m$ coefficient of
$\Li^{-1}g$ is at most $M_m$ times the corresponding maximum for $g$.
Taking $m$th roots and limsups gives \eqref{eq:loss} for finite positive
$\rho(g)$. For $\rho(g)=\infty$, the root growth of $g$ is zero, while
$\limsup M_m^{1/m}=e^\tau<\infty$, so the product again has root growth
zero. The remaining cases and the category assertions follow directly.
\end{proof}

\begin{lemma}[Sharpness and analytic failure]\label{lem:sharp-divisor}
If $\tau>0$ and $0<R<\infty$, there is $\varphi\in\OO(\DR)_2^d$ for
which the formal series $\Li^{-1}\varphi$ is not holomorphic on $\DR$.
It can be chosen to have Taylor radius exactly $Re^{-\tau}$, interpreting
this as zero if $\tau=\infty$.
If $\tau=\infty$, there is even an \emph{entire}
$\varphi\in\OO(\C^d)_2^d$ for which $\Li^{-1}\varphi$ has zero radius.
\end{lemma}
\begin{proof}
Choose increasing degrees $m_k$ along which
$m_k^{-1}\log M_{m_k}\to\tau$, and for each $k$ choose $(\alpha_k,j_k)$
with $|\alpha_k|=m_k$ realizing the maximum inverse divisor. Discard finitely
many degrees so that $M_{m_k}>1$. Then
\begin{equation}\label{eq:finite-radius-witness}
 \varphi(z)=\sum_{k\ge1}R^{-m_k}z^{\alpha_k}e_{j_k}
\end{equation}
is holomorphic on $\DR$. Its inverse has degree-$m_k$ coefficient of
absolute value $R^{-m_k}M_{m_k}$ and no other nonzero coefficients.
Formula \eqref{eq:radius} gives the claimed exact radius.

For the stronger assertion, choose $m_k$ so that
$M_{m_k}\ge\exp(k^2m_k)$, and put
\begin{equation}\label{eq:entire-witness}
 \varphi(z)=\sum_{k\ge1}\exp(-km_k)z^{\alpha_k}e_{j_k}.
\end{equation}
For each ordinary finite $R$, this series converges normally on smaller
closed polydisks once $k$ is large, since $R\exp(-k)\to0$. It is therefore
entire. But its inverse has a degree-$m_k$ coefficient of absolute value at
least $\exp((k^2-k)m_k)$. Its Taylor radius is zero.
\end{proof}
These witnesses use ordinary analytic coefficient functions. No infinitesimal
support manipulation is hidden in their construction.

\section{The nonlinear Hahn linearization theorem}\label{sec:linearization}

\subsection{A linear equation for the unknown change of coordinates}
Let $F=\Lambda z+f$ as in \eqref{eq:Fintro}, and seek $H=z+h$. Define
\begin{equation}\label{eq:Noperator}
 \NN_f h=h(\Lambda z+f(z))-h(\Lambda z)
       =\sum_{|\beta|\ge1}\frac{(\partial^\beta h)(\Lambda z)}{\beta!}f^\beta.
\end{equation}
Schr\"oder's equation becomes
\begin{equation}\label{eq:linear-Schroeder}
 (\Li+\NN_f)h=-f.
\end{equation}
The map $F$ is nonlinear, but with $F$ fixed this equation is \emph{linear in
$h$}. This elementary observation is particularly effective in the Hahn
setting: $\Li$ does not alter Hahn weights, and $\NN_f$ adds positive
support words. Both preserve the absence of coordinate terms of degrees
zero and one.

Whenever $\Li^{-1}$ acts on the coefficient category, let
$T=\Li^{-1}\NN_f$. Formally,
\begin{equation}\label{eq:linearizer-series}
 \boxed{\displaystyle
 h=-\sum_{k\ge0}(-T)^k\Li^{-1}f.}
\end{equation}
The series will now be justified with all support and analytic claims.

\subsection{Sufficiency and uniqueness}
\begin{proposition}[Support-controlled inverse]\label{prop:linearization}
Suppose $B_2^d$ is one of the categories in \cref{thm:main}, and that
$\Li^{-1}$ preserves it. Then \eqref{eq:linearizer-series} defines the unique
normalized positive conjugacy. If $S=\supp f$, its support is contained in
$S^*\setminus\{0\}$. At weight $\gamma$, only terms with
\begin{equation}\label{eq:word-term-bound}
 k+1\le\ell_S(\gamma)
\end{equation}
can contribute, and only input functions $f_\sigma$ with
$\sigma\in D_S(\gamma)$ are involved.
\end{proposition}
\begin{proof}
Each term of $\NN_f$ contains at least one coefficient of $f$ and hence
adds a nonempty word in $S$. The initial input $\Li^{-1}f$ already
contains one support letter. Consequently a term with $k$ copies of $T$
contains at least $k+1$ letters. \Cref{lem:neumann} proves the support
inclusion, strong summability, and \eqref{eq:word-term-bound}. It also says
that only finitely many input support letters and intermediate sums are
relevant at $\gamma$.

Every contributing expression is assembled from finitely many derivatives,
products, compositions with the fixed diagonal rotation, and applications
of $\Li^{-1}$. The assumed category is preserved by those operations. Thus
$h_\gamma$ belongs to $B_2^d$, and \eqref{eq:linearizer-series} is a
legitimate Hahn family in that category. Its operator geometric identity
implies
$(I+T)h=-\Li^{-1}f$, which is \eqref{eq:linear-Schroeder}. Therefore
$H\circ F=\Lambda H$.

For uniqueness, let $h$ and $\widetilde h$ be two normalized positive
solutions, not necessarily with the support just constructed. Their
difference $q$ has well-ordered positive support. If $q\ne0$, take its least
weight $\gamma_0$. Since $\NN_f$ adds positive weights, the coefficient of
$\NN_f q$ at $\gamma_0$ is zero. Equation
$(\Li+\NN_f)q=0$ gives $\Li q_{\gamma_0}=0$. Nonresonance implies
$q_{\gamma_0}=0$, a contradiction. Finally $H$ has a positive inverse on
the same coefficient category by \cref{prop:subgroup}. Its value and
Jacobian at the origin are $0$ and $I$, respectively, by differentiating
the inverse identity. Thus the inverse is normalized as well.
\end{proof}

Applying \cref{lem:radius-loss} gives sufficiency in all five rows of
\cref{thm:main}. Notice that the polynomial row does not require any bound
on the inverse divisors: a polynomial uses only finitely many of them at
any one Hahn coefficient. The infinitely many nonlinear corrections are
separated into different positive Hahn weights.

\subsection{Universal necessity cannot be repaired at later weights}
\begin{proposition}[The first-weight obstruction]\label{prop:first-weight}
Fix $\eta>0$ and an ordinary coefficient vector $\varphi$ vanishing to
order at least two, and let
\[
 F=\Lambda z+t^\eta\varphi(z).
\]
Any normalized positive conjugacy in any of the stated analytic categories
has no nonzero coefficients of weight less than $\eta$, and necessarily
\begin{equation}\label{eq:first-weight}
 h_\eta=-\Li^{-1}\varphi.
\end{equation}
\end{proposition}
\begin{proof}
If a coefficient below $\eta$ is nonzero, take the least nonzero weight of
$h$. At that weight the right side of \eqref{eq:linear-Schroeder} vanishes,
and $\NN_fh$ also vanishes because it adds at least $\eta$. The remaining
equation is $\Li h_\gamma=0$, contradicting nonresonance. At weight
$\eta$, the same reasoning eliminates $\NN_fh$: its input would need
weight at most zero, and $h$ has positive support. The equation is exactly
\eqref{eq:first-weight}.
\end{proof}

\begin{proof}[Completion of the proof of \cref{thm:main}]
For the fixed-domain row, if $\tau>0$ use
\eqref{eq:finite-radius-witness} and set $f=t^\eta\varphi$.
\Cref{prop:first-weight} forces a coefficient which is not holomorphic on
the specified $\DR$. Thus universal preservation of that polydisk fails.

For convergent germs and for entire coefficients, if $\tau=\infty$ use
the entire witness \eqref{eq:entire-witness}. The forced coefficient
$h_\eta$ has zero complex radius, so it is not even a convergent germ.
It certainly cannot be entire. This proves both necessity assertions.
Higher Hahn coefficients, extra scales, and enlargement of $\Gamma$ cannot
change the equation at this first weight. Polynomial and formal sufficiency
was already proved without arithmetic restrictions beyond nonresonance.
The support and uniqueness claims are those of
\cref{prop:linearization}; workspace compatibility is made explicit in
\cref{prop:workspace} below.
\end{proof}

\begin{remark}[Nonresonance is not cosmetic]
If $\lambda^\alpha=\lambda_j$ for some $|\alpha|\ge2$, take
$\varphi=z^\alpha e_j$. Equation \eqref{eq:first-weight} would require
$0=-1$ in that Taylor coefficient. Thus a universal normalized linearization
statement fails in the resonant case even for a polynomial with one
infinitesimal support weight. A different normal form, rather than division
by a zero divisor, is required; \cref{sec:finite-order} supplies it for
finite-order multipliers.
\end{remark}

\section{Quantitative radii, polynomial degrees, and workspace invariance}
\label{sec:quantitative}

\subsection{A finite-depth radius certificate}
\begin{theorem}[Radius at a specified Hahn exponent]\label{thm:radius-depth}
Suppose $\tau<\infty$ and $f=\sum_{\sigma\in S}t^\sigma f_\sigma$
has convergent-germ coefficients. For $\gamma\in S^*\setminus\{0\}$,
choose positive numbers $r_\sigma\le\rho(f_\sigma)$ for the finite set
$\sigma\in D_S(\gamma)$, and put
\[
 r_{\rm in}(\gamma)=\min_{\sigma\in D_S(\gamma)}r_\sigma.
\]
Then the conjugacy in \cref{thm:main} satisfies
\begin{equation}\label{eq:radius-depth}
 \rho(h_\gamma)\ge
       r_{\rm in}(\gamma)\exp\bigl(-\ell_S(\gamma)\tau\bigr).
\end{equation}
If all input coefficients are holomorphic on one $\DR$, this becomes
$\rho(h_\gamma)\ge R\exp(-\ell_S(\gamma)\tau)$.
If all relevant input coefficients are entire, so is $h_\gamma$.
\end{theorem}
\begin{proof}
Consider one term of \eqref{eq:linearizer-series} contributing at $\gamma$.
Its input letters lie in $D_S(\gamma)$. Derivatives, products, and composition
with a diagonal unitary map do not reduce the minimum ordinary Taylor
radius of their inputs. Each application of $\Li^{-1}$ reduces the
available radius by at most $e^{-\tau}$, by \cref{lem:radius-loss}.
A term with $k$ copies of $T$ uses $k+1$ such inversions and has at least
$k+1$ support letters. Hence it uses at most $\ell_S(\gamma)$ inversions.
There are finitely many terms at this weight, so taking their sum preserves
the minimum of these radius bounds. The assertion for entire coefficients
follows either by allowing arbitrarily large finite input radii or directly
from preservation of infinite radius.
\end{proof}
The single-inversion factor is sharp by \cref{lem:sharp-divisor}. We do
not assert that the exponent $\ell_S(\gamma)\tau$ is always attained by a
nonlinear example. Cancellations can improve the radius substantially.
In particular, these lower bounds alone do not decide whether one common
positive radius exists when $0<\tau<\infty$ and word depths are unbounded.

\subsection{A finite algebraic degree certificate}
\begin{theorem}[Polynomial coefficient degrees]\label{thm:degrees}
Suppose every $f_\sigma$ is polynomial and let
$d_\sigma=\max_j\deg(f_{\sigma,j})\ge2$. Then
\begin{equation}\label{eq:degree-bound}
 \deg h_\gamma\le
 1+\max_{(\sigma_1,\ldots,\sigma_k)\in\word_S(\gamma)}
       \sum_{i=1}^k(d_{\sigma_i}-1).
\end{equation}
If all $d_\sigma\le D$, this simplifies to
\begin{equation}\label{eq:uniform-degree}
 \deg h_\gamma\le1+\ell_S(\gamma)(D-1).
\end{equation}
No small-divisor bound is assumed. For polynomial and formal coefficients,
the theorem and the linearization construction also hold for any diagonal
invertible nonresonant $\Lambda$, without $|\lambda_j|=1$.
\end{theorem}
\begin{proof}
The operator $\Li^{-1}$ preserves total polynomial degree. The initial
term associated with one letter $\sigma$ has degree at most $d_\sigma$,
which is $1+(d_\sigma-1)$. Suppose a term associated with a word $w$ has
degree at most $1+\sum_{\sigma\in w}(d_\sigma-1)$, counting repetitions.
A term of $\NN_f$ differentiates it $|\beta|$ times and multiplies by
$|\beta|$ additional coefficient polynomials of $f$. Its degree is at most
\[
 1+\sum_{\sigma\in w}(d_\sigma-1)-|\beta|
        +\sum_{i=1}^{|\beta|}d_{\sigma_i}
 =1+\sum_{\sigma\in w}(d_\sigma-1)
        +\sum_{i=1}^{|\beta|}(d_{\sigma_i}-1).
\]
If the derivative vanishes, there is no term to bound. Induction through
\eqref{eq:linearizer-series}, followed by the finite maximum at $\gamma$,
proves \eqref{eq:degree-bound}. The uniform bound follows. For polynomial
and formal coefficients, composition with any invertible diagonal matrix
preserves the category and the same coefficient equations apply, so the
unit-modulus restriction is unnecessary there.
\end{proof}

\subsection{No new scales and no hidden transfinite input dependency}
\begin{proposition}[Workspace invariance]\label{prop:workspace}
Let $\Gamma\hookrightarrow\Gamma'$ be an ordered group inclusion. The
substitution inverse, positive logarithm, positive exponential, and normalized
nonresonant linearizer constructed above are unchanged when their inputs are
regarded in the larger workspace. Their displayed support certificates
remain valid. In particular, a normalized positive solution constructed in
the larger workspace cannot acquire extra exponents outside $\Gamma$.
Evaluation identities remain valid after adding supports of new halo points.
\end{proposition}
\begin{proof}
The support words use only input exponents and finite addition. The ordinary
coefficient operations, including \eqref{eq:Linverse}, are independent of
$\Gamma$. Thus the explicit formulas compute the same coefficients in
$\Gamma$ and $\Gamma'$. Their support is already inside $\Gamma$.
Uniqueness of substitution inverses, of the positive exponential--logarithm
correspondence, and of normalized linearization identifies these with any
corresponding positive solutions in the larger workspace. Evaluation follows
by including the new positive point supports in the support lemma.
\end{proof}

This does not say that the entire centralizer is unchanged after scalar
extension: a larger workspace can supply new admissible times $c$ in
\eqref{eq:centralizer}. It says that the particular normalized constructions
from the original data remain unchanged.

The finite-dependence assertion is also not an unconditional finite-time
algorithm for arbitrary analytic inputs. It bounds the number of input
\emph{coefficient functions} and analytic operators needed at a specified
Hahn weight. One application of $\Li^{-1}$ to an arbitrary analytic
function can involve infinitely many of its ordinary Taylor coefficients.
With polynomial input and effective finite support certificates, the degree
bound turns each requested Hahn coefficient into a finite exact algebraic
calculation.

\section{Arithmetic separation from ordinary analytic linearization}
\label{sec:arithmetic}

\subsection{The one-variable exponent in continued fractions}
Take $d=1$ and $\lambda=e^{2\pi i\theta}$, where $\theta$ is irrational.
Let $p_k/q_k$ be the convergents of its ordinary continued fraction. The
classical estimates for convergents, recalled for example in
\cite[Appendix A]{CarlettiMarmi}, relate the small-divisor exponent to the
growth of their denominators.

\begin{proposition}[Continued-fraction formula]\label{prop:cf}
For an irrational $\theta$,
\begin{equation}\label{eq:tau-cf}
 \tau(e^{2\pi i\theta})=
 \limsup_{k\to\infty}\frac{\log q_{k+1}}{q_k}.
\end{equation}
\end{proposition}
\begin{proof}
Write $\|x\|$ for distance to the nearest integer. For real $x$,
\[
 4\|x\|\le |e^{2\pi ix}-1|\le2\pi\|x\|.
\]
Consequently the constant factors do not affect the exponential limsup, and
replacing degree $m$ by $n=m-1$ gives
\[
 \tau(\lambda)=\limsup_{n\to\infty}
                  \frac1n\log\frac1{\|n\theta\|}.
\]
The convergent estimates are
\[
 \frac1{q_k+q_{k+1}}<\|q_k\theta\|<\frac1{q_{k+1}},
\]
and best approximation gives
$\|n\theta\|\ge\|q_k\theta\|$ whenever
$q_k\le n<q_{k+1}$. Therefore, apart from a finite initial range,
\[
 \frac{\log q_{k+1}}{q_k}
 <\frac1{q_k}\log\frac1{\|q_k\theta\|}
 \le\frac{\log(q_k+q_{k+1})}{q_k},
\]
and for every $n$ between the two denominators the analogous quotient is
at most the last expression. Its difference from
$(\log q_{k+1})/q_k$ is at most $(\log2)/q_k$, which tends to zero.
Taking limsups proves \eqref{eq:tau-cf}.
\end{proof}

\subsection{A non-Brjuno rotation with unchanged-domain Hahn linearization}
Recall the classical Brjuno condition
\begin{equation}\label{eq:Brjuno}
 \sum_{k\ge0}\frac{\log q_{k+1}}{q_k}<\infty.
\end{equation}
It is strictly stronger than merely requiring the summands to tend to zero.
The latter is exactly $\tau=0$ by \cref{prop:cf}.

\begin{theorem}[An explicit arithmetic separation]\label{thm:nonbrjuno}
There is an explicitly specified irrational rotation $\lambda$ which is not
Brjuno, but for which every positive Hahn perturbation
$F(z)=\lambda z+f(z)$ with coefficient functions holomorphic on a fixed disk
$\mathbb D_R$ and vanishing to order at least two has a normalized positive
linearizer on the \emph{same} coefficient disk.
\end{theorem}
\begin{proof}
Define $\theta=[0;a_1,a_2,\ldots]$ by $a_1=16$ and, recursively for
$k\ge1$,
\begin{equation}\label{eq:cf-construct}
 a_{k+1}=\left\lceil\exp(q_k/k)\right\rceil,\qquad
 q_{k+1}=a_{k+1}q_k+q_{k-1},\qquad q_0=1,\quad q_1=16.
\end{equation}
Every partial quotient is a finite positive integer, so this specifies an
ordinary irrational real number. The denominators grow to infinity. From
\eqref{eq:cf-construct},
\[
 q_k\exp(q_k/k)\le q_{k+1}
      \le q_k\bigl(\exp(q_k/k)+2\bigr).
\]
Taking logarithms and dividing by $q_k$ yields
\[
 \frac1k+\frac{\log q_k}{q_k}
 \le\frac{\log q_{k+1}}{q_k}
 \le\frac1k+\frac{\log q_k}{q_k}
       +\frac{\log(1+2e^{-q_k/k})}{q_k}.
\]
Both outer expressions tend to zero. Hence $\tau(\lambda)=0$ for
$\lambda=e^{2\pi i\theta}$. But the middle expression is at least $1/k$,
so the Brjuno series diverges. The fixed-disk assertion is now
\cref{thm:main}.
\end{proof}

The Brjuno--Yoccoz theorem says that the ordinary quadratic polynomial
$Q_\lambda(w)=\lambda w+w^2$ is locally analytically linearizable precisely
for Brjuno rotations \cite{Yoccoz,BuffCheritat}. Thus the ordinary quadratic
for the rotation just constructed is not analytically linearizable. Our
result does not weaken or contradict that theorem: the ordinary nonlinear
part in $Q_\lambda$ is of size one, whereas the nonlinear part in
\cref{thm:nonbrjuno} has positive Hahn support and no nonzero ordinary
reduction. The resulting conjugacy is a Hahn family, not necessarily a
convergent ordinary family under a numerical specialization of $t$.

\begin{example}[Every threshold regime occurs]
For any prescribed $a>0$, replace \eqref{eq:cf-construct} by
$a_{k+1}=\lceil\exp(aq_k)\rceil$. The same estimates give
$(\log q_{k+1})/q_k\to a$, so $\tau=a$.
Taking $a_{k+1}=\lceil\exp(kq_k)\rceil$ gives $\tau=\infty$.
Therefore the finite positive and infinite rows of the classification are
not empty arithmetic possibilities.
\end{example}

\subsection{An entire-coefficient conjugacy that stops at a critical scale}
\begin{proposition}[Quadratic scaling and the boundary of the halo theorem]
\label{prop:quadratic-scale}
For every non-root-of-unity $\lambda\in\C^*$ and every $\eta>0$,
\begin{equation}\label{eq:scaled-quadratic}
 F(z)=\lambda z+t^\eta z^2
\end{equation}
has a normalized conjugacy with polynomial Hahn coefficients
\begin{equation}\label{eq:scaled-H}
 H(z)=z+\sum_{n\ge1}b_n t^{n\eta}z^{n+1}.
\end{equation}
It evaluates on every finite halo, irrespective of the growth of $b_n$.
For every nonzero ordinary $w$, substituting $z=t^{-\eta}w$ into the
individual terms of \eqref{eq:scaled-H} is \emph{not} a strongly summable
family. If $\lambda$ is a non-Brjuno irrational rotation, the ordinary series
$w+\sum_{n\ge1}b_nw^{n+1}$ has zero radius of convergence as well.
\end{proposition}
\begin{proof}
Let
\[
 L_\lambda(w)=w+\sum_{n\ge1}b_nw^{n+1}
\]
be the unique formal solution of
$L_\lambda(\lambda w+w^2)=\lambda L_\lambda(w)$; nonresonance gives its
coefficients recursively. Substitute $w=t^\eta z$ and divide the equation
by $t^\eta$. The resulting expression is exactly
\eqref{eq:scaled-H} and satisfies $H\circ F=\lambda H$. Each Hahn
coefficient is polynomial. The support is $\{n\eta:n\ge1\}$ or a subset,
which is well ordered. Evaluation at a finite point is justified by the
support lemma, with no condition on the ordinary magnitudes of the $b_n$.

Infinitely many $b_n$ are nonzero. Otherwise $L_\lambda$ would be a
nonconstant polynomial, and the degrees of
$L_\lambda\circ Q_\lambda$ and $\lambda L_\lambda$ would be respectively
$2\deg L_\lambda$ and $\deg L_\lambda$, an impossibility. At the indicated
critical scale every nonzero summand becomes
\[
 b_n t^{n\eta}(t^{-\eta}w)^{n+1}
      =t^{-\eta}b_nw^{n+1}.
\]
Thus infinitely many members have the same Hahn exponent $-\eta$, violating
\cref{def:summable}. This is a failure of local finiteness, not of a
real-valued norm estimate.

Finally, if the formal series $L_\lambda$ had a positive ordinary radius,
it would be a local holomorphic coordinate with derivative one and would
analytically linearize $Q_\lambda$. For a non-Brjuno rotation this contradicts
the classical theorem just cited. Hence its radius is zero.
\end{proof}

Even for a Brjuno rotation, ordinary convergence of the coefficient series
would not by itself make the repeated-weight family a strong Hahn sum.
It would be an additional ordinary analytic summation convention. The
non-Brjuno example rules out even that local analytic reinterpretation of
the formal linearizer. We make no assertion about sectorial or resurgent
summation under additional hypotheses.

\subsection{Arithmetic invisible to the Hahn valuation}
For every nonzero ordinary divisor $\lambda^\alpha-\lambda_j$, its Hahn
valuation is zero. A proof keeping only the Hahn valuation therefore cannot
distinguish the small-divisor regimes above. The distinction comes from the
\emph{ordinary complex Taylor growth inside each coefficient function}.
This explains why formal coefficients always permit normalized
linearization under nonresonance, yet analytic coefficient categories have
sharp arithmetic thresholds.

In particular, the negative results do not prohibit every conjugacy on an
infinitesimal monad. The formal-coefficient conjugacy still exists and can be
evaluated there. They prohibit a conjugacy in the specified, smaller
analytic coefficient category. This distinction is indispensable when
comparing different definitions of ``surcomplex analytic''.

\section{Finite-order multipliers and resonant rotation--flow normal forms}
\label{sec:finite-order}

\subsection{A finite averaging construction}
Nonresonance is false for finite-order multipliers. Nevertheless there is an
exact normal form which keeps precisely the resonant dynamics, without any
small-divisor estimates.

\begin{theorem}[Same-domain finite-order decomposition]\label{thm:finite-order}
Let $\Lambda$ be diagonal unitary with $\Lambda^q=I$ for an integer $q\ge1$,
and let $F=\Lambda z+f$, where
$f\in\HH_\Gamma(\OO(\DR))_+^d$ vanishes coefficientwise to order at least
two. There exist a normalized positive coordinate change $H$ and a positive
vector field $Y$ vanishing to order at least two, both on the same $\DR$,
such that
\begin{equation}\label{eq:resonant-normal}
 H\circ F\circ H^{-1}=\Lambda\circ\Phi_Y,
 \qquad Y(\Lambda z)=\Lambda Y(z).
\end{equation}
The rotation and flow in \eqref{eq:resonant-normal} commute. Every monomial
$z^\alpha e_j$ of $Y$ satisfies the resonance condition
$\lambda^\alpha=\lambda_j$. All coefficients have finite support-word
certificates from $\supp f$, and the construction commutes with workspace
enlargement. Moreover,
\begin{equation}\label{eq:finiteorder-equivalence}
 F^{\circ q}=\id
 \quad\Longleftrightarrow\quad Y=0
 \quad\Longleftrightarrow\quad
 F\text{ is positively conjugate to }\Lambda.
\end{equation}
The same statement holds for germs, entire coefficients, polynomials, and
formal coefficients.
\end{theorem}
\begin{proof}
The reduction of $P=F^{\circ q}$ is identity. It fixes zero and has exact
Jacobian $I$ there. By \cref{thm:exp-log},
\[
 X=\log C_P
\]
is a positive vector field whose coefficients vanish to order at least two.
The map $F$ commutes with $P$, so its pullback commutes with $\log C_P$.
Equivalently, $F$ preserves $X$ and commutes with all its admissible ordinary
flow times. Set
\begin{equation}\label{eq:Afinite}
 A=F\circ\Phi_{-X/q}.
\end{equation}
The factors commute, and hence
$A^{\circ q}=P\circ\Phi_{-X}=\id$. Its ordinary reduction is $\Lambda$.
Define the finite average
\begin{equation}\label{eq:average}
 H(z)=\frac1q\sum_{j=0}^{q-1}\Lambda^{-j}A^{\circ j}(z).
\end{equation}
The leading term of every summand after multiplication by $\Lambda^{-j}$
is $z$, so $H=\id+\text{positive}$, with value zero and derivative $I$ at
zero. It has a positive inverse on the same coefficient domain. Reindexing
the finite sum, using $A^q=\id$ and $\Lambda^q=I$, gives
\[
 H\circ A=\Lambda H.
\]

For clarity, the pushforward vector field is defined in coordinates by
\[
 (H_*X)(w)=dH(H^{-1}(w))\,a(H^{-1}(w)),\qquad X=a\cdot\partial.
\]
Taylor substitution shows that it is positive and has the same analytic
coefficient category. It conjugates flows:
$H\circ\Phi_X\circ H^{-1}=\Phi_{H_*X}$, either by the coefficientwise
chain rule or by conjugating pullback operators. Put
\[
 Y=\frac1qH_*X.
\]
Since $A$ preserves $X$, the conjugated map $\Lambda$ preserves $Y$.
This is exactly $Y(\Lambda z)=\Lambda Y(z)$, which also gives the claimed
monomial resonance condition. From $F=A\circ\Phi_{X/q}$ we obtain
\eqref{eq:resonant-normal} and its commuting-factor assertion.

The operations used are finite compositions with the ordinary rotation,
positive logarithms, exponentials, substitution inverses, and the finite
sum \eqref{eq:average}. Each preserves the original coefficient domain and
the generated support monoid, by earlier results. Finally, $P=\id$ if and
only if $X=0$, which holds if and only if $Y=0$. If $Y=0$, the displayed
normal form is a conjugacy to $\Lambda$; conversely such a conjugacy forces
$F^q=\id$. This proves \eqref{eq:finiteorder-equivalence}.
\end{proof}

\subsection{Interpretation and limits of the normal form}
In one variable with $\lambda$ of exact order $q$, the resonant vector
field has the form
\[
 Y(z)=\sum_{\substack{n\ge2\\n\equiv1\pmod q}}a_n z^n\partial_z,
\]
interpreted coefficientwise when $a_n$ themselves are Hahn series.
The coefficient family, rather than an unrestricted rearrangement in $n$,
remains the defining object. If $q=1$, the theorem reduces to the positive
flow correspondence.

The finite-factor/flow splitting follows the classical
semisimple--unipotent pattern \cite{Ribon}. We do not claim that the finite
averaging idea or that abstract pattern is new. The point established here
is that the entire construction is valid for arbitrary positive Hahn support
and retains the same ordinary coefficient domain. Nor do we claim uniqueness
of every resonant normal form: further coordinate changes commuting with
$\Lambda$ can transform $Y$. Formula \eqref{eq:average} specifies the
particular normalized coordinate used in this construction.

\section{Worked examples and exact coefficient calculations}
\label{sec:examples}

\subsection{The iterative logarithm of a quadratic perturbation}
Take $\Gamma=\Q$, write $t=t^1$, and consider $F(z)=z+tz^2$ with entire
polynomial coefficients. Directly from \eqref{eq:explog},
\begin{align}\label{eq:quadratic-log}
 X_F={}&\biggl(tz^2-t^2z^3+\frac32t^3z^4-\frac83t^4z^5
                +\frac{31}{6}t^5z^6-\frac{157}{15}t^6z^7
                +O(t^7)\biggr)\partial_z.
\end{align}
Here $O(t^7)$ refers only to this single-parameter formal expansion.
Exponentiating half of this vector field produces the unique positive
half-iterate of $F$, with polynomial coefficients at every Hahn weight.
The ordinary size and growth of the rational coefficients in
\eqref{eq:quadratic-log} have no bearing on strong summability in $t$.

One can see the origin of the first terms for a general map. Introduce a
bookkeeping parameter $\varepsilon$ and take $F_\varepsilon=z+\varepsilon u$.
Expanding the logarithm through degree three in that parameter gives
\begin{align}\label{eq:general-third-log}
 X_{F_\varepsilon}z={}&\varepsilon u
 -\frac{\varepsilon^2}{2}(Du)u\\
 &+\varepsilon^3\left(\frac1{12}D^2u(u,u)
                    +\frac13(Du)^2u\right)+O(\varepsilon^4).\notag
\end{align}
Indeed, $\Delta^2z=\varepsilon^2(Du)u+
\frac12\varepsilon^3D^2u(u,u)+O(\varepsilon^4)$, while
$\Delta^3z=\varepsilon^3(D^2u(u,u)+(Du)^2u)+O(\varepsilon^4)$.
Substitution in the logarithmic series proves \eqref{eq:general-third-log}.
With $u=z^2$ it gives the first three terms of \eqref{eq:quadratic-log}.

\subsection{A complete centralizer and a double fixed point}
For $X=tz^2\partial_z$, the flow has the especially simple expression
\begin{equation}\label{eq:mobius-flow}
 \Phi_c(z)=\frac{z}{1-ctz}
          =z+\sum_{n\ge1}(ct)^n z^{n+1},\qquad
 c=0\text{ or }v(c)>-1.
\end{equation}
The last series is a strongly summable expansion on finite halos. The time-one
map in this example is $z/(1-tz)$, not the quadratic map from the previous
subsection. By \cref{thm:centralizer}, the full positive centralizer of this
time-one map on the entire-coefficient halo consists exactly of
\eqref{eq:mobius-flow}. For every nonzero admissible time,
\[
 \Phi_c(z)-z=\frac{ctz^2}{1-ctz}.
\]
The factor $ct/(1-ctz)$ is a unit in the Hahn coefficient algebra. Thus the
fixed-point ideal is exactly $(z^2)$ for every such time: the point at zero
retains its scheme-theoretic double structure, not merely its location.

\subsection{Quadratic linearization coefficients}
For the normalized ordinary formal equation
$L_\lambda(\lambda w+w^2)=\lambda L_\lambda(w)$, the first coefficients
in \cref{prop:quadratic-scale} are
\begin{align}\label{eq:quadratic-bs}
 b_1&=-\frac1{\lambda(\lambda-1)},\\
 b_2&=\frac2{\lambda(\lambda-1)^2(\lambda+1)},\notag\\
 b_3&=-\frac{5\lambda^2+1}
 {\lambda^2(\lambda-1)^3(\lambda+1)(\lambda^2+\lambda+1)}.\notag
\end{align}
They are obtained by equating successive powers of $w$; the denominators
at stage $n$ arise from $\lambda^{n+1}-\lambda$. The fourth coefficient
and exact symbolic verification of all four are included in
\texttt{data/verification.json}. Multiplying $b_nz^{n+1}$ by $t^{n\eta}$
turns these ordinary formal data into the strongly summable surcomplex
linearizer, even when the ordinary series has zero radius.

\subsection{Two independent perturbation scales}
Let $t$ and $u$ represent $t^1$ and $t^\Omega$ in
$\Gamma=\Q+\Q\Omega$. Consider the map with multiplier $-1$,
\[
 F(z)=-z+tz^2+uz^3.
\]
It illustrates \cref{thm:finite-order} with $q=2$. The averaging coordinate
and the resonant vector field begin as
\begin{align}\label{eq:resonant-example}
 H(z)&=z-\frac{t}{2}z^2+\frac{t^2}{2}z^3-\frac{5tu}{4}z^4
                  +O((t,u)^3),\\
 Y(z)&=\left(-uz^3-t^2z^3-\frac32u^2z^5
              -\frac{25}{4}t^2uz^5-\frac72u^3z^7
              +O((t,u)^4)\right)\partial_z.\notag
\end{align}
In these displays, $O((t,u)^N)$ means discarded monomials of \emph{total
parameter degree} at least $N$ in an auxiliary two-parameter jet algebra.
It is not a valuation cutoff: for example $t^{100}$ has much smaller Hahn
valuation than $u$, although its total parameter degree is larger.
Every displayed monomial of $Y$ has odd coordinate degree, exactly as the
resonance condition $Y(-z)=-Y(z)$ requires. The normal form satisfies
$HF=(-\Phi_Y)H$.

The verification program also uses the two-dimensional identity-reducing map
\[
 G(x,y)=(x+tx^2+uy,\ y+txy+ux).
\]
At total parameter depth three it checks the logarithm, inverse, exponential,
Leibniz identity, and the fixed-point factorization
\eqref{eq:fixedfactor}. This example includes infinitesimal linear terms,
which are permitted in the identity-reducing flow theorem even though they
are excluded from the exact-multiplier linearization theorem.

\subsection{An arbitrary-support example beyond a finite grid}
Use $S$ and $\Gamma$ from \cref{ex:ranktwo}, and let
$F=\Lambda z+\sum_{\sigma\in S}t^\sigma f_\sigma(z)$, with each
$f_\sigma$ vanishing to order at least two. Under a valid analytic
hypothesis from \cref{thm:main}, the higher-rank coefficients obey
\begin{align}\label{eq:Omega-example}
 h_\Omega&=-\Li^{-1}f_\Omega,\\
 h_{\Omega+1}&=\Li^{-1}\left(
       D(\Li^{-1}f_\Omega)(\Lambda z)f_1(z)
      +D(\Li^{-1}f_1)(\Lambda z)f_\Omega(z)\right).\notag
\end{align}
Only the word $(\Omega)$ contributes in the first line, and only the words
$(\Omega,1)$ and $(1,\Omega)$ can contribute in the second. Infinitely
many input coefficients at rational weights below $\Omega$ do not enter
either formula. This is a direct use of finite word dependence, rather than
a rank-one approximation to the value group.

\section{Literature comparison and the scope of the novelty claim}
\label{sec:literature}

\subsection{The directly relevant predecessors}
The positive support calculus is classical Hahn--Neumann mathematics.
\Cref{lem:neumann} is explicitly attributed to that theory, with Higman's
lemma furnishing the proof we use \cite{Higman,Poonen}. The support lemmas
are not claimed as discoveries of this paper.

The formal exponential--logarithm correspondence and its role in embedding
maps in vector-field flows are also classical. Rib\'on's treatment gives
formal logarithms and uses the Jordan--Chevalley decomposition in precisely
this general circle of ideas \cite{Ribon}. Our arguments recast that
mechanism with a different filtration, by arbitrary positive Hahn support,
and prove preservation of an entire family of ordinary coefficient domains.
The centralizer and finite-order splitting should be understood as
support-and-domain versions of familiar formal dynamical structures.

The ordinary homological equation, the denominators
$\lambda^\alpha-\lambda_j$, and the relation between analytic linearization
and arithmetic are standard. Carletti--Marmi study linearization in analytic
and nonanalytic classes and are a close predecessor for the distinction
between a formal solution and regularity of that solution
\cite{CarlettiMarmi}. The Brjuno--Yoccoz quadratic theorem is used as an
external classical theorem, not reproved here \cite{Yoccoz,BuffCheritat}.
Our contribution is a different universal classification: coefficient
regularity is demanded separately at each Hahn exponent, and finite
support depth determines how many homological inversions affect one such
coefficient.

Berarducci--Mantova interpret transseries and larger classes of series as
surreal analytic functions and study composition and differentiation
\cite{BM}. Those functions and their derivatives should not be conflated
with the present coordinate-holomorphic coefficient algebra. Likewise,
Cluckers--Lipshitz develop a broad theory of valued fields with analytic
structure \cite{CL}. We do not assert that the coefficient categories here
lie outside every framework encompassed by that work. Theorems must be
compared with their coefficient rings, growth requirements, and evaluation
domains stated, rather than by the word ``analytic'' alone.

\subsection{Comparison with the motivating repository}
The repository catalogue, the analysis and analytic-geometry READMEs, and
the opening framework of the analytic-geometry article were reviewed at the
snapshot specified on the title page \cite{Repo}.
They already distinguish fixed domains, common-domain germs, radius-free
germs, and formal coefficients, and contain extensive local geometry and
residue theory. The present manuscript does not repeat that program as a
new contribution. It develops dynamics of substitutions, sharp homological
regularity thresholds, iteration-invariant fixed-point ideals, centralizers,
and resonant splittings. A targeted repository search for
\texttt{linearization} returned no indexed matches; that search alone is not
a proof of absence from every file. The catalogue and the reviewed
foundational documents supplied the substantive basis for selecting this
direction.

None of the repository's unrefereed Noetherianity, Nullstellensatz, residue,
or global geometry claims is used as a theorem in these proofs. The needed
support statements and coordinate operations are established here from the
classical inputs already identified. In particular,
\cref{thm:fixed-ideal} compares ideals by an explicit invertible matrix;
it does not infer ideal equality from equality of zero sets.

\subsection{Precisely delimited proposed contributions}
The strongest proposed new statement is \cref{thm:main}, together with its
finite-depth estimates and first-weight obstructions. Its universal
quantifiers, its distinctions among the five coefficient categories, and
its absence of a grid or rank restriction are part of the statement.
\Cref{thm:nonbrjuno,prop:quadratic-scale} explain a concrete separation from
ordinary analytic dynamics and the sharp need for a stated halo domain.
The same-domain versions of \cref{thm:exp-log,thm:fixed-ideal,thm:centralizer,thm:finite-order}
form a supporting dynamical package.

A targeted search of these themes was conducted on September 21, 2026,
including formal flow embedding, Hahn-supported analytic functions,
small-divisor linearization, and the named repository. It found the substantial
predecessors above, but no exact published statement of the combined
arbitrary-support classification proved here. This is a limited, auditable
novelty assessment, not a guarantee of first priority. The paper proves
candidate new theorems rather than claiming to have settled a named
previously published conjecture. Independent expert review of both the
proofs and the literature comparison remains appropriate.

\section{Remaining questions and non-claims}\label{sec:limitations}

\subsection{A genuine gap between the two analytic categories}
\begin{question}[Common-domain germs for finite positive small-divisor exponent]
\label{q:common-germ}
Assume $0<\tau(\Lambda)<\infty$. Characterize the positive perturbations
with coefficients holomorphic on one common ordinary polydisk for which the
unique radius-free linearizer has coefficients holomorphic on some common
possibly smaller polydisk.
\end{question}
Our fixed-disk counterexample shows that the original polydisk may be lost.
Our coefficientwise estimates guarantee a positive radius for each fixed
Hahn coefficient. Neither statement decides whether the infimum of all
these radii is positive. An exact characterization would need to control
interactions across unbounded word depths rather than a single homological
inverse. The lower bound in \eqref{eq:radius-depth} is not a proof that
radius collapse actually occurs.

\subsection{Resonances of infinite order and higher-dimensional centralizers}
\begin{question}[General resonant reduction]
For a diagonal unitary multiplier with resonances but without finite order,
find a support-controlled analytic normal form with precise necessary and
sufficient coefficient-category hypotheses.
\end{question}
Nonresonant monomials can be treated by the homological inverse, but resonant
terms create a coupled normal-form problem and remove the simple global
inverse used in \eqref{eq:linearizer-series}. A finite cyclic averaging
operator is no longer available. We have not proved a universal analytic
Poincar\'e--Dulac statement in that setting.

In several variables, classification of the positive centralizer reduces
exactly to classification of the commuting positive vector fields. The
one-variable meromorphic ratio argument cannot classify those higher-
dimensional Lie centralizers. Additional first integrals or commuting
geometric structures can occur.

\subsection{No all-scale, topological, or unrestricted summation assertion}
The evaluated common-domain results concern finite halos; germ and formal
coefficient results concern the infinitesimal monad. The critical-scale
example \cref{prop:quadratic-scale} explicitly prevents automatic extension
to all infinite surcomplex arguments. We do not construct a global analytic
atlas on all of $\No[i]$, an ordinary contour theory on its fine topology,
or a fine-topologically continuous complex-time flow.

There is no assertion that every algebraic conjugacy of evaluated functions
must arise from a coefficient family, or that failure in a specified analytic
category prevents all set-theoretic conjugacies. Theorems are identities in
stated Hahn coefficient algebras and hence identities of their evaluated
maps. No interchange of an infinite ordinary coefficient sum with a
non-locally-finite Hahn family is permitted.

\section{Computational checks and reproducibility}\label{sec:verification}
The accompanying program \texttt{code/verify.py} implements a sparse exact
polynomial jet algebra. Coefficients lie in $\Q(i)$, with a symbolic
$\lambda$ used for four displayed ordinary linearizer coefficients. The
truncation is by total degree in one or two independent perturbation
parameters. All comparisons are exact; no floating-point tolerance is used.

The recorded run passes \textbf{21 checks}. They include exponential after
logarithm, two-sided inversion, half and third iterates, time addition,
Leibniz identities, identification of the logarithm with its coordinate
vector field, the two-variable fixed-point matrix factorization, normalized
nonresonant conjugacy, its inverse and degree bound, and the finite-order
averaging and resonant flow identities. The one-parameter flow checks run
through degree six. The two-parameter multivariable, nonresonant, and
resonant checks run through total parameter degree three.

Run the check suite from this directory with
\begin{verbatim}
python code/verify.py
\end{verbatim}
The full list of assertions and example expansions is recorded in
\texttt{data/verification.json}; the console transcript is in
\texttt{data/verification.txt}. The script requires Python 3.10 or later and
SymPy. The article builds without an external bibliography or graphics file:
\begin{verbatim}
latexmk -pdf -interaction=nonstopmode article.tex
\end{verbatim}

These checks validate the finite algebraic formulas used as illustrations.
They are not an implementation of arbitrary Hahn fields, not a proof of the
well-order or analytic radius assertions, and not a proof-assistant
verification. No Lean verification is claimed. In particular, truncating by
total parameter degree in \cref{sec:examples} does not approximate the
entire ordered-value-group filtration. The general theorems rest on the
proofs in the text, especially the finite-word certificate and the
first-weight obstruction.

\appendix
\section{A short proof of the word-order input}\label{app:higman}
We record the special case of Higman's lemma used in
\cref{lem:neumann}, so that the support proof can be audited without a
separate combinatorial argument being hidden in the notation. It is a
classical result, credited in the main text \cite{Higman}.

Let $S$ be well ordered. A word $(a_1,\ldots,a_m)$ embeds in a word
$(b_1,\ldots,b_n)$ if there are indices
$i_1<\cdots<i_m$ with $a_j\le b_{i_j}$ for every $j$. A sequence of words
is called bad if no earlier word embeds in a later one. Suppose an infinite
bad sequence exists. Choose a bad sequence $w_0,w_1,\ldots$ minimally by
word length at each stage: after choosing a prefix, choose a shortest next
word among those which admit a bad infinite extension. The choices take
place in a set; ordinary choice suffices. No $w_n$ is empty.

Write $w_n=v_n a_n$, separating the last letter. Any infinite sequence in a
well-ordered set has an infinite nondecreasing subsequence. One direct way
to choose one is to select a position attaining the minimum of the current
infinite tail, then repeat on the tail strictly after that position. The
successive minima cannot decrease. Thus choose
$n_0<n_1<\cdots$ with $a_{n_0}\le a_{n_1}\le\cdots$.

Now consider
\[
 w_0,\ldots,w_{n_0-1},v_{n_0},v_{n_1},v_{n_2},\ldots.
\]
This is bad. An earlier $w_i$ embedding in a later $v_{n_j}$ would also
embed in $w_{n_j}$, contradicting the original badness. If
$v_{n_i}$ embeds in $v_{n_j}$ for $i<j$, append the dominated last letters
to obtain an embedding of $w_{n_i}$ in $w_{n_j}$, again a contradiction.
Comparisons within the unchanged prefix were already forbidden. The new
bad sequence agrees with the chosen one until position $n_0$ and has the
strictly shorter word $v_{n_0}$ there. This contradicts the minimal choice.
No bad sequence exists, which is the required word well-quasi-order.

When $S\subset\Gamma_{>0}$, a proper word embedding strictly raises the
sum: either an unused positive letter is present in the larger word, or a
matched letter is strictly larger. This last observation converts the
word theorem into both well-ordering of $S^*$ and finiteness of the words of
one weight, as claimed in \cref{lem:neumann}.

\section{A coefficientwise recipe for a certified calculation}
\label{app:recipe}
Suppose one wishes to compute $h_\gamma$ for polynomial input in
\cref{thm:main}, and has an effective presentation of the finite word set
$\word_S(\gamma)$. The following recipe specifies what must be certified.

First collect the finite letter set $D_S(\gamma)$ and the finite set of
partial sums of all words of weight $\gamma$. Only input coefficient
polynomials at these letters can matter. The maximum word length gives a
bound $L=\ell_S(\gamma)$. In \eqref{eq:linearizer-series}, retain only
$k=0,\ldots,L-1$, and in each Taylor expansion of $\NN_f$ retain only
products whose support words can complete an allowed word of total weight
$\gamma$. Polynomial degrees can be bounded in advance by
\eqref{eq:degree-bound}.

Apply $\Li^{-1}$ to each retained monomial by the exact scalar division in
\eqref{eq:Linverse}. Nonresonance of the finitely many divisors actually
used must be established, or supplied as an exact symbolic condition. Sum
the resulting finite expressions and check Schr\"oder's equation at the
requested weight. The word certificate proves that discarded expressions
cannot affect it. This is a finite exact algorithm \emph{relative to the
support certificate and the polynomial coefficient data}.

For analytic input, the same finite support certificate still applies, but
$\Li^{-1}$ is an operator on an entire coefficient function, not merely a
finite list of Taylor coefficients. A radius certificate from
\cref{thm:radius-depth} is then part of the analytic output. A finite jet
calculation alone must not be described as computing the entire analytic
coefficient, nor as proving its convergence.

\section{Notation and hypotheses at a glance}\label{app:notation}
\begin{longtable}{@{}P{0.23\textwidth}P{0.70\textwidth}@{}}
\toprule
Symbol or condition & Meaning and role\\
\midrule
\endhead
$\Gamma$ & Nonzero set-sized ordered abelian group; take a divisible
subgroup of $\No$ for the surcomplex normal-form interpretation.\\[0.35em]
$K_\Gamma$ & Hahn field $\C((t^\Gamma))$, with $v(t^\gamma)=\gamma$.\\[0.35em]
$\HH_\Gamma(B)$ & Hahn families with coefficients in the specified ordinary
algebra $B$; no uniform ordinary norm bound is included.\\[0.35em]
$+$ & All nonzero support exponents are strictly positive. It does not mean
that ordinary complex coefficients are positive.\\[0.35em]
$U_\Gamma^\#$ & Finite halo $U+\m_{K_\Gamma}^d$, the evaluated domain for
ordinary-domain coefficient families.\\[0.35em]
$B_2$ & Coefficients vanishing at zero together with their first derivatives.
Used for exact-multiplier normalized linearization.\\[0.35em]
$S^*$ & Finite sums of positive letters from $S$, including zero.\\[0.35em]
$\ell_S(\gamma)$ & Maximum length of a support word of total weight
$\gamma$; finite even when infinitely many support elements precede
$\gamma$.\\[0.35em]
$\Li$ & Homological operator $h(\Lambda z)-\Lambda h$ on coordinate terms
of degree at least two.\\[0.35em]
$\tau(\Lambda)$ & Exponential rate of the inverse ordinary complex
homological divisors, not a Hahn valuation.\\[0.35em]
$\rho(g)$ & Isotropic ordinary Taylor radius of one coefficient function.\\[0.35em]
$X_F$ & Positive coordinate vector field $\log C_F$ when $F$ reduces to
identity. All Hahn scalars are constant for its coordinate derivatives.\\[0.35em]
$\Phi_X$ & Time-one map $(\exp X)z$; for a time $c$, use $\Phi_{cX}$.\\[0.35em]
$I_{\Fix(F)}$ & Ideal generated by $F_j-z_j$, compared by an explicit
invertible matrix, not merely by its zero set.\\
\bottomrule
\end{longtable}

\begin{thebibliography}{99}
\raggedright

\bibitem{Repo}
Vladimir Reshetnikov,
\emph{Surreal}, research and formalization repository,
\texttt{docs/README.md},
\texttt{docs/surcomplex/analysis/README.md}, and
\texttt{docs/surcomplex/analytic-geometry/}.
Snapshot commit
\texttt{39f2be6667ade51bca2b45daa47e289d69c09764},
reviewed September 21, 2026. Unrefereed AI-assisted research drafts.
\url{https://github.com/VladimirReshetnikov/Surreal}.

\bibitem{Higman}
G. Higman,
Ordering by divisibility in abstract algebras,
\emph{Proceedings of the London Mathematical Society} (3) \textbf{2}
(1952), 326--336.
\href{https://doi.org/10.1112/plms/s3-2.1.326}{doi:10.1112/plms/s3-2.1.326}.

\bibitem{Poonen}
B. Poonen,
\emph{Units in Hahn--Mal'cev--Neumann rings},
manuscript dated January 14, 2026, 4 pp.
\url{https://math.mit.edu/~poonen/papers/malcev.pdf}.

\bibitem{Ribon}
J. Rib\'on,
Embedding smooth and formal diffeomorphisms through the Jordan--Chevalley
decomposition,
\emph{Journal of Differential Equations} \textbf{253} (2012), no.~12,
3211--3231.
\href{https://doi.org/10.1016/j.jde.2012.09.009}{doi:10.1016/j.jde.2012.09.009}.
Preprint: \href{https://arxiv.org/abs/1107.3601}{arXiv:1107.3601}.

\bibitem{CarlettiMarmi}
T. Carletti and S. Marmi,
Linearization of analytic and non-analytic germs of diffeomorphisms of
$(\C,0)$,
\emph{Bulletin de la Soci\'et\'e Math\'ematique de France}
\textbf{128} (2000), no.~1, 69--85.
\href{https://doi.org/10.24033/bsmf.2363}{doi:10.24033/bsmf.2363}.
\url{https://www.numdam.org/articles/10.24033/bsmf.2363/}.

\bibitem{Yoccoz}
J.-C. Yoccoz,
Th\'eor\`eme de Siegel, nombres de Bruno et polyn\^omes quadratiques,
in \emph{Petits diviseurs en dimension 1},
\emph{Ast\'erisque} \textbf{231} (1995), 3--88.
Publisher's volume page:
\url{https://smf.emath.fr/publications/petits-diviseurs-en-dimension-1}.

\bibitem{BuffCheritat}
X. Buff and A. Ch\'eritat,
\emph{The Brjuno function continuously estimates the size of quadratic
Siegel disks}, 2004 preprint,
\href{https://arxiv.org/abs/math/0401044}{arXiv:math/0401044}.

\bibitem{BM}
A. Berarducci and V. Mantova,
Transseries as germs of surreal functions,
\emph{Transactions of the American Mathematical Society}
\textbf{371} (2019), 3549--3592.
\href{https://doi.org/10.1090/tran/7428}{doi:10.1090/tran/7428}.
\href{https://arxiv.org/abs/1703.01995}{arXiv:1703.01995}.

\bibitem{CL}
R. Cluckers and L. Lipshitz,
Fields with analytic structure,
\emph{Journal of the European Mathematical Society}
\textbf{13} (2011), no.~4, 1147--1223.
\href{https://doi.org/10.4171/JEMS/278}{doi:10.4171/JEMS/278}.

\end{thebibliography}
\end{document}
